% ============================================================================
\clearpage
\section*{v4.9.0-r13 Source-Only: The Time Scale of the Kicked Duality Dynamics}
\addcontentsline{toc}{section}{v4.9.0-r13 Source-Only: The Time Scale of the Kicked Duality Dynamics}
\label{sec:v490-r13-time-scale}
% ============================================================================

\paragraph{Purpose.}
The dynamics proposed in r11 fixes only the product \(\omega\tau=\pi/6\).  This layer derives the time scale from the clock
model the source already contains (Section~\ref{sec:v470-r15-zero-clock-phase-window-relay-stitching}).  It follows the
source in reading the photon-clock probe as a diagnostic timing function, not as an identity between Riemann zeros and
photons.  All statements are local, in the clock's own units, and make no global photon-clock statement.  The script is
\texttt{llm/k54\_zeitskala.py}.  Earlier layers are not rewritten.

\noindent\textbf{Source-only status.}
This v4.9.0-r13 layer is source-only.  It does not build a PDF, does not assign a new DOI,
and does not authorize a public release.  The v4.9.0 Version DOI
\href{https://doi.org/10.5281/zenodo.23534954}{10.5281/zenodo.23534954} and the Concept DOI
\href{https://doi.org/10.5281/zenodo.19210728}{10.5281/zenodo.19210728} are unchanged by it.

\subsection*{v4.9.0-r13 the rate from the tick profile}
\label{subsec:v490-r13-rate}

r15 writes the local clock tick as \(G_V(\tau)=1+(G_V^{\rm avail}-1)\cos^2\bigl((\tau-\tau_V)/\sigma_V\bigr)\).  Between two kicks of
r11 the energy in the compensator port is \(\cos^2(\omega t)\).  Both are \(\cos^2\) of an angle linear in time.

\paragraph{Theoremlet v4.9.0-r13.1.}
\begin{enumerate}[leftmargin=2em]
\item[(a)] Identifying the two profiles gives \(\omega=1/\sigma_V\): the tick width of r15 is the time the duality rotation of r11
needs for one radian.  The kick interval is \(\tau_{\rm k}=\tfrac\pi6\,\sigma_V\).
\item[(b)] The smallest probe with a profile has two terms.  For
\(\mathcal C_\gamma\propto|e^{i\gamma_a\tau}+e^{i\gamma_b\tau}|^2\) the coherent part is exactly
\(\cos^2\bigl(\tfrac12(\gamma_b-\gamma_a)\tau\bigr)\), so \(\sigma=2/(\gamma_b-\gamma_a)\).  The beat period is
\(T_{\rm b}=2\pi/(\gamma_b-\gamma_a)\).
\item[(c)] Hence \(\tau_{\rm k}=T_{\rm b}/6\) for every pair of clock frequencies.  Three stages, the quarter turn from matter to
tachyon, take half a beat.  Six stages, the half turn \(z\mapsto-z\), take one beat.  Twelve take two beats, and twenty-four,
the return of the spinor, take four.  The detector stage \(n=20\) lies at \(10/3\) beats.
\item[(d)] For the two lowest frequencies the source uses, \(\gamma_1=14.1347\ldots\) and \(\gamma_2=21.0220\ldots\):
\(\sigma=0.29039\), \(T_{\rm b}=0.91228\), \(\tau_{\rm k}=0.15205\) in clock units.
\end{enumerate}

\paragraph{Proof sketch.}
(b) \(|e^{ia}+e^{ib}|^2=4\cos^2\tfrac{b-a}2\); checked numerically to \(10^{-15}\).  (c) \(\tau_{\rm k}=\tfrac\pi6\cdot\tfrac2{\Delta\gamma}=
\tfrac16\cdot\tfrac{2\pi}{\Delta\gamma}\).  \(\square\)

\paragraph{What is fixed and what is not.}
The relative statement (c) holds for any pair of frequencies.  The absolute value (d) depends on which pair carries the
tick, and the source names none.  The first two are the minimal choice.  The clock unit itself has no conversion to
seconds in the source.  Supplying one would require identifying clock frequencies with physical frequencies.  The source
expressly declines that identification, and as a global statement it would touch the forbidden global photon-clock
statement.  The time scale is therefore derived up to the unit of the clock.

\subsection*{v4.9.0-r13 the windows of the primary chain}
\label{subsec:v490-r13-windows}

\paragraph{Theoremlet v4.9.0-r13.2.}
In the angle of r13.1(a), the half-widths \(\delta_V/\sigma_V\) of r15 are \(39.23^\circ\) for \(V3\), \(50.77^\circ\) for \(V5\)
and \(28.13^\circ\) for \(V8\).  A timing offset of one C5 stage, \(30^\circ\), stays inside every window of the primary chain
\(\mathcal R_A=\{V3,V5\}\) and leaves the window of \(V8\) on the chain \(\mathcal R_B\).

\paragraph{Proof sketch.}
Conversion of the r15 table.  \(\square\)

\paragraph{Reading.}
This is a second, independent reason for the restriction (P4) of r12 to \(\mathcal R_A\): only there does a single C5 stage
fit into the clock window.  The tolerance \(26.57^\circ\) of r12 concerns the loss of clock content per readout, a different
quantity from a timing offset within one window, so the two numbers do not contradict each other.  The window of \(V8\)
lies between them.

\subsection*{v4.9.0-r13 the stage angle between Hopf and lattice}
\label{subsec:v490-r13-hopf-lattice}

The main body relates the Hopf value and the measured value of the offset by
\(\bsym=\tfrac1{\sqrt3}\exp\bigl(-\gamma_1/(4\pi\gL^5)\bigr)=0.5492\) (the table with ``Hopf (continuum \(S^3\))'' and
``Physical (prime lattice)'' rows).  The stage angle \(\arctan\bsym\) of r10 is therefore \(30^\circ\) in the continuum limit
\(\gamma_1\to0\) and \(28.78^\circ\) on the lattice, with \(q=0.7683\).  The C5 factor \(3/4\) is the continuum value.  On the
lattice the Hopf reading would give stage \(22\) (r10), while the triolic reading keeps \(\pi/6\).  Which of the two describes
the C5 recursion is open.  The detector stage \(n=20\) favours the continuum value or the triolic reading.

\subsection*{v4.9.0-r13 register entries}
\label{subsec:v490-r13-register}

\begin{center}
\small
\begin{tabular}{@{}>{\raggedright\arraybackslash}p{0.34\textwidth}>{\raggedright\arraybackslash}p{0.58\textwidth}@{}}
\toprule
statement & status after v4.9.0-r13\\
\midrule
time scale of r11 & \(\omega=1/\sigma_V\), \(\tau_{\rm k}=T_{\rm b}/6\) (one sixth of the clock beat) for every frequency pair (r13.1)\\
absolute value & \(\tau_{\rm k}=0.15205\) clock units for the two lowest frequencies; no conversion to seconds, by the source's own restriction\\
r12 premise (P4) & second reason: one C5 stage fits only into the windows of \(\mathcal R_A\) (r13.2)\\
C5 factor \(3/4\) & continuum value of the Hopf reading; lattice value \(0.7683\); triolic reading unaffected\\
\bottomrule
\end{tabular}
\end{center}

\subsection*{v4.9.0-r13 ledger and guard update}
\label{subsec:v490-r13-ledger-guard}

\[
R_{\rm time}^{v490\text{-}r13}=(0,0,1)
\]
for (relative time scale derived, windows of the primary chain checked, unit of the clock open).

The layer spends no guard:
\[
R_{\rm glob}^{(v490\text{-}r13,+)}=(0,1,1,1,1),
\qquad
\Gamma_{\rm guard}^{v490\text{-}r13}=(1,1,1,1,1),
\qquad
\beta_4=1/256\ \text{unspent}.
\]
It does not imply exact finite \(C5\) zero closure, exact all-readout terminality,
\(C6/T2\), a global photon-clock/RH theorem, or global all-readout closure.

